%! TEX program = lualatex
%! TeX encoding = UTF-8 Unicode
%! TeX spellcheck = fr-FR
%! BIB TS-program = biber
% -*- coding: UTF-8; -*-
%------------------------------------


%------------------------------------
\newcommand{\Title}{Curriculum Vit\ae}
\newcommand{\Author}{Laurent Perrinet}
\newcommand{\Address}{Researcher in Computational Neuroscience (DR2 CNRS)\\ Institut de Neurosciences de la Timone\\
UMR 7289, CNRS / Aix-Marseille Universit{\'e}\\ 27, Bd. Jean Moulin, 13385 Marseille Cedex 5, France}%
\newcommand{\Website}{https://laurentperrinet.github.io}
\newcommand{\Email}{Laurent.Perrinet@univ-amu.fr}
%------------------------------------
\documentclass[11pt, a4paper]{article}
\usepackage{fontspec}
% DOCUMENT LAYOUT
\usepackage{geometry}
\geometry{a4paper, textwidth=13cm, textheight=25cm, marginparsep=7pt, marginparwidth=3cm}
\setlength\parindent{0in}
% FONTS
\usepackage[usenames,dvipsnames]{color}
\usepackage{xltxtra}
\defaultfontfeatures{Mapping=tex-text}
% ---- CUSTOM COMMANDS
\chardef\&="E050
\newcommand{\html}[1]{\href{#1}{\scriptsize\textsc{[html]}}}
\newcommand{\pdf}[1]{\href{#1}{\scriptsize\textsc{[pdf]}}}
\newcommand{\doi}[1]{\href{#1}{\scriptsize\textsc{[doi]}}}
\newcommand{\amper}{{\selectfont\itshape\&}}
% ---- MARGIN YEARS
\usepackage{marginnote}
\newcommand{\years}[1]{\marginnote{\scriptsize #1}}
\renewcommand*{\raggedleftmarginnote}{}
\setlength{\marginparsep}{7pt}
\reversemarginpar
% HEADINGS
\usepackage{sectsty}
\usepackage[normalem]{ulem}
\sectionfont{\sffamily\mdseries\large\underline}
\subsectionfont{\rmfamily\mdseries\scshape\normalsize}
\subsubsectionfont{\rmfamily\bfseries\upshape\normalsize}
% PDF SETUP
\usepackage[bookmarks, colorlinks, breaklinks]{hyperref}
\hypersetup{linkcolor=blue,citecolor=blue,filecolor=black,urlcolor=blue,
pdftitle={\Title},%
pdfauthor={\Author < \Email > \Address - \Website},%
pdfsubject={\Title}%
}%
\expandafter\ifx\csname urlstyle\endcsname\relax
  \providecommand{\doi}[1]{doi: #1}\else
  \providecommand{\doi}{doi: \begingroup \urlstyle{rm}\Url}\fi

% DOCUMENT
\begin{document}
\pagestyle{empty}
\reversemarginpar
\textsf{\Large \Author}\\[1cm]
\Address \\[.2cm]
\textsc{url}: \href{\Website}{\Website}\\

\section*{Research interests}
My research program aims to elucidate the relationship between neural structure and sensory function: how do the topographic organization of neurons and the spiking nature of neural signals enable sensory systems to optimally adapt to the statistical regularities of natural scenes, through predictive processing mechanisms? This work opens the way to novel neuromorphic algorithms with high energy efficiency.

\section*{Areas of specialization}
Spatio-temporal inference in low-level sensory areas.\\
Unsupervised learning in topographic maps.\\
Predictive coding, sparse representations and active perception.

\section*{Education}

\noindent\years{2017}\textsc{Habilitation à Diriger des Recherches} (HDR), Aix-Marseille Université. Title: ``Predictive coding in visuo-motor transformations''.

\noindent\years{1999--2003}\textsc{PhD} in Cognitive Science, ONERA/DTIM \& Université Paul Sabatier, Toulouse (France). Highest distinction with jury congratulations.

\noindent\years{1993--1998}\textsc{MSc} in Engineering, {\sc Supaéro} (Toulouse, France). Specialization in stochastic models for signal and image processing.

\section*{Professional experience}

\noindent\years{2020--\ldots}\textsc{Research Director} (DR2, CID51), Institut de Neurosciences de la Timone, CNRS / Aix-Marseille Université, Marseille.

\noindent\years{2004--2020}\textsc{Research Scientist} (CR2 then CR1), INCM then INT, CNRS, Marseille.

\noindent\years{2010--2012}\textsc{Long-term visit}, Karl Friston's theoretical neurobiology group, Wellcome Trust Centre for Neuroimaging, University College London, UK.

\noindent\years{2004}\textsc{Post-doctoral fellow}, Redwood Neuroscience Institute, San Francisco, USA. Supervisor: Bruno Olshausen.

\section*{Supervision}

\noindent 6 PhD theses directed (completed), 4 co-directed (completed), 5 current PhD students (3 as principal supervisor, 2 co-supervised), 5 post-docs supervised. Reviewer for 18 PhD defenses.

\section*{Selected grants}

\noindent\years{2023--2027} \textit{Emergences} -- Near-physics emerging models for embedded AI (PEPR France 2030).

\noindent\years{2021--2025} PI: ANR \textit{AgileNeuRobot} (ANR-20-CE23-0021) -- Agile bio-mimetic aerial robots for real-world flight, 435\,k€.

\noindent\years{2022--2025} PI: \textit{Polychronies} (A*MIDEX AMX-21-RID-025) -- Spatio-temporal spike motifs, 100\,k€.

\noindent\years{2021--2026} Collaborator: ANR \textit{PRIOSENS} (2021--2025), ANR \textit{ACES} (2022--2026), ANR \textit{RubinVase} (2021--2024).

\section*{Selected publications}

\noindent\years{2025} Andrew Isaac Meso, Jonathan Vacher, Nikos Gekas, Pascal Mamassian, Laurent U Perrinet, Guillaume S Masson. ``DynTex: A Real-Time Generative Model of Dynamic Naturalistic Luminance Textures.'' {\bf Journal of Vision}.

\noindent\years{2025} Skye Gunasekaran \emph{et al.}, Laurent Perrinet. ``A Predictive Approach to Enhance Time-Series Forecasting.'' {\bf Nature Communications}.

\noindent\years{2024} Antoine Grimaldi, Victor Boutin, Sio-Hoi Ieng, Ryad Benosman, Laurent U Perrinet. ``A Robust Event-Driven Approach to Always-on Object Recognition.'' {\bf Neural Networks}.

\noindent\years{2023} Hugo Ladret, Nelson Cortes, Lamyae Ikan, Frédéric Chavane, Christian Casanova, Laurent U Perrinet. ``Cortical recurrence supports resilience to sensory variance in the primary visual cortex.'' {\bf Nature Communications Biology}.

\noindent\years{2022} Frédéric Chavane, Laurent U Perrinet, James Rankin. ``Revisiting Horizontal Connectivity Rules in V1: From like-to-like towards like-to-All.'' {\bf Brain Structure and Function}.

\noindent\years{2021} Victor Boutin, Angelo Franciosini, Franck Ruffier, Frédéric Chavane, Laurent U Perrinet. ``Sparse Deep Predictive Coding captures contour integration capabilities of the early visual system.'' {\bf PLoS Computational Biology}.

\noindent\years{2020} Chloé Pasturel, Anna Montagnini, Laurent U Perrinet. ``Humans adapt their anticipatory eye movements to the volatility of visual motion properties.'' {\bf PLoS Computational Biology}.

\noindent\years{2014} Laurent U Perrinet, Rick A Adams, Karl Friston. ``Active inference, eye movements and oculomotor delays.'' {\bf Biological Cybernetics}.

\noindent\years{2012} Karl Friston, Rick A Adams, Laurent U Perrinet, Michael Breakspear. ``Perceptions as Hypotheses: Saccades as Experiments.'' {\bf Frontiers in Psychology}.

\noindent\years{2012} Claudio Simoncini, Laurent U Perrinet, Anna Montagnini, Pascal Mamassian, Guillaume S Masson. ``More is not always better: dissociation between perception and action explained by adaptive gain control.'' {\bf Nature Neuroscience}.

\noindent\years{2010} Laurent Perrinet. ``Role of homeostasis in learning sparse representations.'' {\bf Neural Computation}.

\noindent\years{2004} Laurent Perrinet, Manuel Samuelides, Simon Thorpe. ``Coding static natural images using spiking event times: do neurons cooperate?'' {\bf IEEE Transactions on Neural Networks}.

\medskip
\noindent{\footnotesize Bibliometric indices (March 2026): 63 peer-reviewed journal articles, $>$4200 citations, h$=$30, i10$=$61. \href{https://scholar.google.com/citations?user=TVyUV38AAAAJ}{Google Scholar}.}

{\footnotesize \href{https://github.com/laurentperrinet/perrinet_curriculum-vitae.tex}{Last updated: \today} --- \href{https://github.com/laurentperrinet/perrinet_curriculum-vitae.tex/blob/master/perrinet_curriculum-vitae-full.pdf}{Full version}.}

\end{document}